% ================================================================
\subsubsection{MPC Architecture}
\label{sec:mpc_arch}
% ================================================================

\paragraph{Block-diagram overview.}
Figure~\ref{fig:mpc_block} shows the closed-loop architecture.
At each timestep $k$ the observer produces the estimated augmented state
$\hat{\tilde{\vxi}}_k = [\hat{\vxi}_k^\top,\hat{W}_k]^\top$
(identical machinery to Section~\ref{sec:lqr_arch}: leaky kinematic integrator,
AoA direct measurement, $C_L$ bisection for gust estimation, and optional EKF
correction).
The MPC block uses $\hat{\tilde{\vxi}}_k$ as the initial condition for a
horizon-$N$ optimisation over the control sequence
$\mathbf{u} = [\delta_k, \delta_{k+1}, \ldots, \delta_{k+N-1}]^\top$,
solves it via Adam with TensorFlow automatic differentiation, and applies only
the first element $\delta_k^*$ to the plant (receding-horizon principle).

\begin{figure}[htbp]
\centering
\begin{tikzpicture}[>=latex,
    block/.style={draw, rectangle, minimum height=1.0cm, minimum width=2.8cm,
                  align=center, font=\small},
    tallblock/.style={draw, rectangle, minimum height=1.8cm, minimum width=2.8cm,
                  align=center, font=\small},
    lbl/.style={font=\small}]

  % Online blocks
  \node[block]                              (plant)    {Plant\\(LDNet + struct.)};
  \node[block, below=2.0cm of plant]        (observer) {Observer};
  \node[tallblock, left=3.4cm of plant]     (mpc)      {MPC solver\\(Adam, $N$ steps)\\$\min_{\mathbf{u}}\;J_N$};

  % plant -> observer (right side, measurements)
  \draw[->] (plant.east) -- ++(1.0,0) coordinate (mr)
            -- node[right, lbl]{$\ddot{h},\;\alpha,\;C_L$}
            (mr |- observer.east) -- (observer.east);

  % observer -> mpc (bottom, estimated state)
  \draw[->] (observer.west) -- ++(-0.8,0) coordinate (sl)
            -- node[below, lbl]{$\hat{\tilde{\vxi}}_k = [\hat{h},\hat{\dot{h}},\hat{\alpha},\hat{\dot{\alpha}},\hat{z},\hat{W}]^\top$}
            (sl |- mpc.west) -- (mpc.west);

  % mpc -> plant (top, first element of optimal sequence)
  \draw[->] (mpc.east) -- node[above, lbl]{$\delta_k^* = \mathbf{u}^*[0]$} (plant.west);

  % receding-horizon annotation
  \node[font=\footnotesize, text=gray, below=0.15cm of mpc]
        (rh) {(shift \& warm-start for $k+1$)};

\end{tikzpicture}
\caption{MPC closed-loop architecture.
         The observer is identical to the LQR case (Section~\ref{sec:lqr_arch}).
         At each timestep the MPC solver optimises the full $N$-step control
         sequence over the nonlinear rollout, then applies only the first
         element $\delta_k^*$ to the plant (receding-horizon principle).}
\label{fig:mpc_block}
\end{figure}

\paragraph{Receding-horizon principle.}
Although the solver produces a full sequence
$\mathbf{u}^* = [\delta_k^*, \ldots, \delta_{k+N-1}^*]^\top$,
only $\delta_k^*$ is applied.
At the next timestep $k+1$, the optimisation is repeated from the updated
state estimate $\hat{\tilde{\vxi}}_{k+1}$, with the remaining elements
$[\delta_{k+1}^*, \ldots, \delta_{k+N-1}^*, 0]$ shifted by one and used as
the warm-start for the Adam solver (Section~\ref{sec:warm_start}).
This closed-loop re-planning makes MPC implicitly robust to model mismatch and
observer error: any deviation from the predicted trajectory is corrected at the
next solve.

\paragraph{Gust sequence over the horizon.}
\label{sec:gust_seq}
The plant propagator requires a gust input $W_j$ at each horizon step $j$.
In the absence of a dedicated gust sensor, this is approximated from the
observer estimate $\hat{W}_k$ and its one-step finite-difference rate:
\begin{equation}
  \hat{W}_{k+j} \approx \hat{W}_k + j\,\Delta t\,\frac{\hat{W}_k - \hat{W}_{k-1}}{\Delta t}
  = \hat{W}_k + j\,(\hat{W}_k - \hat{W}_{k-1}),
  \quad j = 0,\ldots,N-1,
  \label{eq:W_extrap}
\end{equation}
clipped to $[0, 80]\,\text{m/s}$.
This linear extrapolation is crude but gives the optimizer a physically
consistent disturbance preview even without a dedicated LiDAR or upwash sensor.
In the oracle (\emph{true-state}) configuration the exact future gust profile
$W_{k}, \ldots, W_{k+N-1}$ is passed directly to the solver, providing an
upper bound on achievable performance.

% ================================================================
\subsubsection{Problem Formulation}
\label{sec:mpc_formulation}
% ================================================================

\paragraph{State and augmented dynamics.}
The MPC operates on the same augmented state
$\tilde{\vxi}_k = [\vxi_k^\top, W_k]^\top \in \mathbb{R}^6$
as the LQR (Section~\ref{sec:lqr}), where
$\vxi_k = [h_k, \dot{h}_k, \alpha_k, \dot{\alpha}_k, z_k]^\top \in \mathbb{R}^5$.
Unlike the LQR, however, the MPC does \emph{not} linearise the dynamics:
it propagates the full nonlinear map
\begin{equation}
  \vxi_{k+1} = F_{\rm nl}(\vxi_k, \delta_k, W_k)
  \label{eq:mpc_propagator}
\end{equation}
defined by the LDNet aerodynamic step \eqref{eq:zdyn}--\eqref{eq:yrec}
followed by the structural RK4 integrator \eqref{eq:rhs}.
The gust evolves by the same exponential decay assumed in the observer:
$W_{k+1} = \lambda_w W_k$, $\lambda_w = 0.98$.

\paragraph{Finite-horizon cost.}
Given the current state estimate $\hat{\tilde{\vxi}}_k$, the MPC solves at
each timestep
\begin{equation}
    \begin{split}
        \min_{\mathbf{u} \in \mathcal{U}^N} J_N(\mathbf{u};\,\hat{\tilde{\vxi}}_k) = \sum_{j=0}^{N-1} \Bigl[ & Q_{C_L} C_{L,k+j}^2 + Q_{C_M} C_{M,k+j}^2 + Q_h\, h_{k+j}^2 + Q_\alpha\, \alpha_{k+j}^2 \\ 
        & + Q_{\Delta C_L}(\Delta C_{L,k+j})^2 + R\,\delta_{k+j}^2 + R_{\Delta u}(\delta_{k+j} - \delta_{k+j-1})^2 \Bigr],
    \end{split}
    \label{eq:mpc_cost}
\end{equation}
where $\Delta C_{L,k+j} = C_{L,k+j} - C_{L,k+j-1}$ is the one-step lift
increment, and the predicted states are obtained by rolling out
\eqref{eq:mpc_propagator} from $\hat{\tilde{\vxi}}_k$ under the sequence
$\mathbf{u} = [\delta_k, \ldots, \delta_{k+N-1}]^\top$.

\paragraph{Constraint set.}
The only hard constraint is the actuator saturation:
\begin{equation}
  \mathcal{U} = \{\delta \in \mathbb{R} : |\delta| \le \delta_{\max} = 20^\circ\}.
  \label{eq:box}
\end{equation}
No rate constraint is imposed in the cost above; the term $R_{\Delta u}$ acts
as a soft rate penalty that discourages abrupt changes without rendering the
problem infeasible.

\paragraph{Weight selection.}
The primary tuning objective is GLA: suppress $C_L$ peaks and structural
oscillations under a $60\,\text{m/s}$ gust.
The baseline weights used throughout this chapter are:
\begin{equation}
  Q_{C_L} = 10, \quad
  Q_{C_M} = 1, \quad
  Q_h = \frac{1}{0.00428^2} \approx 5.5 \times 10^4, \quad
  Q_\alpha = \frac{5}{0.00823^2} \approx 7.4 \times 10^4,
  \label{eq:mpc_weights_state}
\end{equation}
\begin{equation}
  Q_{\Delta C_L} = 5, \quad
  R = 1, \quad
  R_{\Delta u} = 0.1,
  \label{eq:mpc_weights_ctrl}
\end{equation}
with horizon $N = 10$ steps ($0.1\,\text{s}$ preview) and Adam learning rate
$\eta = 0.5$ for $n_{\rm iter} = 15$ iterations per timestep.
The $Q_h$ and $Q_\alpha$ normalisations match those used in the LQR
\eqref{eq:weights}; the $Q_{C_L}$ and $Q_{\Delta C_L}$ terms have no LQR
counterpart and directly penalise aerodynamic load excursions, including their
rate of change---a capability that is nontrivial to encode in a static Riccati
weight.
Sensitivity to these choices is discussed in Section~\ref{sec:mpc_sensitivity}.

% ================================================================
\subsubsection{Gradient-Based Solver: Backpropagation Through the Rollout}
\label{sec:mpc_solver}
% ================================================================

The cost \eqref{eq:mpc_cost} is a scalar function of $\mathbf{u} \in \mathbb{R}^N$
defined through the composition of $N$ nonlinear plant steps.
Because the plant passes through the LDNet weights, no closed-form gradient
$\nabla_{\mathbf{u}} J_N$ exists; it is instead computed via reverse-mode AD,
exactly as the Jacobian matrices were computed for the LQR design
(Section~\ref{sec:jacobian})---but now applied to the \emph{full horizon} rather
than a single step, and repeated online at every timestep.

\paragraph{The rollout as a computation graph.}
Define the $j$-th rollout step as the map
\begin{equation}
  \Psi_j : (\vxi_j, \delta_j, \hat{W}_{k+j}) \;\longmapsto\;
  \bigl(\vxi_{j+1},\; C_{L,j},\; C_{M,j}\bigr),
  \label{eq:Psi_j}
\end{equation}
evaluated with the predicted gust input $\hat{W}_{k+j}$.
Internally, $\Psi_j$ is a finite composition of two types of elementary
differentiable operations, applied in sequence:
\begin{enumerate}
  \item \textbf{Affine layers} of LDNet.  Denoting by $\mathbf{a}^{(\ell)}$
        the vector of activations (outputs) at layer $\ell$, each affine layer
        computes
        $\mathbf{a}^{(\ell)} = \mathbf{W}^{(\ell)}\mathbf{a}^{(\ell-1)} + \mathbf{b}^{(\ell)}$,
        where $\mathbf{W}^{(\ell)}$ and $\mathbf{b}^{(\ell)}$ are the frozen
        weight matrix and bias vector of layer $\ell$, and
        $\mathbf{a}^{(0)}$ is the network input
        $[z,\, h,\, \dot{h},\, \alpha,\, \dot{\alpha},\, \delta,\, W]^\top$.
        The Jacobian of this operation with respect to its input $\mathbf{a}^{(\ell-1)}$
        is simply $\mathbf{W}^{(\ell)}$---a constant, already in memory.
  \item \textbf{Element-wise activations.}  Each activation layer applies a
        nonlinearity $\sigma$ component by component:
        $\mathbf{a}^{(\ell)} \leftarrow \sigma(\mathbf{a}^{(\ell)})$.
        Its Jacobian is the diagonal matrix \\
        $\operatorname{diag}(\sigma'(a^{(\ell)}_1), \ldots, \sigma'(a^{(\ell)}_n))$,
        where $\sigma'$ is the scalar derivative of the activation function
        (e.g.\ $\sigma'(x) = 1 - \tanh^2(x)$ for a $\tanh$ unit).
\end{enumerate}
These two types cover all operations in both LDNet sub-networks
($f_{\rm dyn}$ and $f_{\rm rec}$).
After the aerodynamic step, $\Psi_j$ applies the structural RK4 integrator,
which computes the four slope estimates
\begin{equation}
    \begin{aligned}
        k_1 &= g(x_j,\, F_y,\, M_z), \quad & k_2 &= g\!\left(x_j + \tfrac{\Delta t}{2}k_1,\, F_y,\, M_z\right), \\
        k_3 &= g\!\left(x_j + \tfrac{\Delta t}{2}k_2,\, F_y,\, M_z\right), \quad & k_4 &= g(x_j + \Delta t\,k_3,\, F_y,\, M_z),
    \end{aligned}
    \label{eq:rk4_slopes}
\end{equation}
and then updates the structural state as
$x_{j+1} = x_j + \tfrac{\Delta t}{6}(k_1 + 2k_2 + 2k_3 + k_4)$.
The right-hand side $g$ in \eqref{eq:rhs} is \emph{linear} in the structural
state $x = [h, \dot{h}, \alpha, \dot{\alpha}]^\top$: it consists of linear
combinations of the state components plus the aerodynamic forces $F_y$ and
$M_z$, which arrive from LDNet and are held constant during the RK4 sub-steps.
Each slope $k_i$ is therefore an affine function of $x_j$---exactly the same
type of primitive as an affine layer in the network.
From the perspective of automatic differentiation, RK4 requires no special
treatment: TensorFlow records each sub-step as an additional affine operation
in the computation graph and differentiates through all four stages
automatically, in exactly the same way it differentiates through the network
layers.
Composing $N$ such steps gives the rollout map
\begin{equation}
  \mathcal{R}(\mathbf{u};\,\hat{\tilde{\vxi}}_k)
  = \Psi_{N-1} \circ \cdots \circ \Psi_1 \circ \Psi_0
  \bigl(\hat{\vxi}_k,\; \mathbf{u},\; \hat{\mathbf{W}}\bigr),
  \label{eq:rollout}
\end{equation}
where $\hat{\mathbf{W}} = [\hat{W}_{k}, \ldots, \hat{W}_{k+N-1}]^\top$ is the
gust sequence \eqref{eq:W_extrap}.
Since every elementary operation in every $\Psi_j$ is differentiable, the full
rollout $\mathcal{R}$ is differentiable almost everywhere with respect to
$\mathbf{u}$, and its gradient can be computed exactly by applying the chain
rule backwards through the entire composition---with no finite differences and
no approximation of the network weights.

\paragraph{Gradient derivation via the adjoint method.}
\label{sec:adjoint}
Let $\ell_j(C_{L,j}, C_{M,j}, h_j, \alpha_j, \delta_j, \delta_{j-1})$ denote
the per-step cost in \eqref{eq:mpc_cost}, so that
$J_N = \sum_{j=0}^{N-1} \ell_j$.
Differentiating $J_N$ with respect to the control element $\delta_m$
($0 \le m \le N-1$) by the chain rule gives
\begin{equation}
  \frac{\partial J_N}{\partial \delta_m}
  = \underbrace{\sum_{j=m}^{N-1}
    \frac{\partial \ell_j}{\partial \vxi_j}\,
    \frac{\partial \vxi_j}{\partial \delta_m}}_{\text{indirect: }\delta_m \to \vxi_{m+1} \to \cdots \to \ell_j}
    + \underbrace{\frac{\partial \ell_m}{\partial \delta_m}}_{\text{direct: }\delta_m \to \ell_m},
  \label{eq:dJ_ddelta}
\end{equation}
where the two terms capture two distinct channels through which $\delta_m$
affects the total cost.
The \emph{indirect} channel is the effect of $\delta_m$ on all future costs
$\ell_j$ ($j \ge m$) via the state trajectory: $\delta_m$ drives $\vxi_{m+1}$
through the plant, which drives $\vxi_{m+2}$, and so on, altering every
subsequent cost term.
The \emph{direct} channel is the immediate penalty that $\delta_m$ incurs in
$\ell_m$ regardless of the state, through the control-effort terms
$R\,\delta_m^2 + R_{\Delta u}(\delta_m - \delta_{m-1})^2$:
these depend on $\delta_m$ explicitly and do not pass through any state
variable, so they are not captured by the sum over $\partial\ell_j/\partial\vxi_j$.
The sensitivity $\partial \vxi_j / \partial \delta_m$ in the indirect term
propagates forward through the composition $\Psi_{j-1} \circ \cdots \circ \Psi_m$.
\emph{Forward-mode} differentiation computes this sensitivity by seeding a
unit perturbation $\dot{\delta}_m = 1$ at the input and propagating it forward
through the chain step by step, accumulating $\partial \vxi_j / \partial \delta_m$
as it goes.
This must be repeated once for each control element $\delta_m$, $m = 0,\ldots,N-1$,
and for each $m$ the perturbation must be carried through all $j \geq m$ steps:
the total cost is $O(N^2)$ forward passes.
The fundamental inefficiency is that forward mode computes one \emph{column}
of the Jacobian per pass; with $N$ inputs and a scalar output $J_N$, this is
the wrong direction.

\emph{Reverse-mode} AD (backpropagation) exploits the scalar output to reverse
this trade-off: because $J_N \in \mathbb{R}$ is a single number, one backward
pass propagates its gradient with respect to \emph{all} $N$ inputs
simultaneously.
The general rule is: forward mode is efficient when there are few inputs and
many outputs; reverse mode is efficient when there are many inputs and few
outputs---exactly the case here ($N$ control inputs, one scalar cost).
Reverse-mode AD reduces the total cost to a \emph{single} forward pass to
record activations, followed by a \emph{single} backward pass, regardless of $N$.6
Define the adjoint (co-state) vector at step $j$ as
\begin{equation}
  \boldsymbol{\lambda}_j = \frac{\partial J_N}{\partial \vxi_j}
  \in \mathbb{R}^5,
  \label{eq:costate}
\end{equation}
which satisfies the backward recursion
\begin{equation}
  \boldsymbol{\lambda}_{j-1}
  = \left(\frac{\partial \Psi_{j-1}}{\partial \vxi_{j-1}}\right)^\top \boldsymbol{\lambda}_j
    + \left(\frac{\partial \ell_{j-1}}{\partial \vxi_{j-1}}\right)^\top,
  \qquad
  \boldsymbol{\lambda}_{N} = \mathbf{0},
  \label{eq:adjoint_rec}
\end{equation}
with terminal condition $\boldsymbol{\lambda}_N = \mathbf{0}$ (no terminal cost).
Each Jacobian--vector product $(\partial \Psi_j / \partial \vxi_j)^\top \boldsymbol{\lambda}_{j+1}$
is itself computed by one backward pass through the primitive chain
$\phi_1, \ldots, \phi_K$ of step $j$, exactly as in \eqref{eq:adjoint}.

The gradient with respect to $\delta_m$ is then recovered as the
$m$-th partial of the loss through the control coupling:
\begin{equation}
  \frac{\partial J_N}{\partial \delta_m}
  = \left(\frac{\partial \Psi_m}{\partial \delta_m}\right)^\top \boldsymbol{\lambda}_{m+1}
    + \frac{\partial \ell_m}{\partial \delta_m},
  \label{eq:grad_u}
\end{equation}
so the full gradient $\nabla_{\mathbf{u}} J_N \in \mathbb{R}^N$ is assembled in
$O(N)$ backward passes, one per rollout step, with no additional computational
graph traversals.
In practice TensorFlow's \texttt{GradientTape} records the forward computation
automatically during the rollout and evaluates \eqref{eq:adjoint_rec}--\eqref{eq:grad_u}
in a single \texttt{tape.gradient(J, u\_var)} call.


%%%%%%%%%%%%%%%%%%%%
Ok, proviamo con un esempio piccolo: $N = 3$ passi, quindi $\mathbf{u} = [\delta_0, \delta_1, \delta_2]$.

Il problema: vuoi $\partial J_N / \partial \delta_0$, $\partial J_N / \partial \delta_1$, $\partial J_N / \partial \delta_2$.

La catena causale è:


δ₀ → ξ₁ → ξ₂ → ξ₃
δ₁ →      ξ₂ → ξ₃
δ₂ →           ξ₃
Ogni $\ell_j$ dipende da $\vxi_j$ e da $\delta_j$ direttamente.

L'idea dell'adjoint: invece di chiedersi "quanto cambia $J_N$ se perturbo $\delta_m$?" (forward, $N$ calcoli), ti chiedi "quanto vale $J_N$, e poi distribuisco la responsabilità all'indietro".

Parti dalla fine. Nessun costo terminale, quindi:
$$\boldsymbol{\lambda}_3 = \mathbf{0}$$

Ora vai indietro di un passo. $\vxi_3$ ha contribuito a $\ell_2$ (l'unico costo che dipende da $\vxi_3$... aspetta, $\ell_j$ dipende da $\vxi_j$ non $\vxi_{j+1}$). Più precisamente:

$$\boldsymbol{\lambda}_2 = \left(\frac{\partial \Psi_2}{\partial \vxi_2}\right)^\top \boldsymbol{\lambda}_3 + \left(\frac{\partial \ell_2}{\partial \vxi_2}\right)^\top = \mathbf{0} + \left(\frac{\partial \ell_2}{\partial \vxi_2}\right)^\top$$

$$\boldsymbol{\lambda}_1 = \left(\frac{\partial \Psi_1}{\partial \vxi_1}\right)^\top \boldsymbol{\lambda}_2 + \left(\frac{\partial \ell_1}{\partial \vxi_1}\right)^\top$$

$$\boldsymbol{\lambda}_0 = \left(\frac{\partial \Psi_0}{\partial \vxi_0}\right)^\top \boldsymbol{\lambda}_1 + \left(\frac{\partial \ell_0}{\partial \vxi_0}\right)^\top$$

Cos'è fisicamente $\boldsymbol{\lambda}_j$? È "se cambio $\vxi_j$ di un po', di quanto cambia il costo totale futuro $J_N$?". È la sensibilità del costo totale rispetto allo stato al passo $j$ — una sorta di "prezzo" dello stato in quel momento.

Perché questo è efficiente? Una volta calcolati $\boldsymbol{\lambda}_0, \ldots, \boldsymbol{\lambda}_3$ con una sola passata all'indietro, i gradienti rispetto a tutti i controlli vengono fuori gratis:

$$\frac{\partial J_N}{\partial \delta_2} = \left(\frac{\partial \Psi_2}{\partial \delta_2}\right)^\top \boldsymbol{\lambda}_3 + \frac{\partial \ell_2}{\partial \delta_2}$$

$$\frac{\partial J_N}{\partial \delta_1} = \left(\frac{\partial \Psi_1}{\partial \delta_1}\right)^\top \boldsymbol{\lambda}_2 + \frac{\partial \ell_1}{\partial \delta_1}$$

$$\frac{\partial J_N}{\partial \delta_0} = \left(\frac{\partial \Psi_0}{\partial \delta_0}\right)^\top \boldsymbol{\lambda}_1 + \frac{\partial \ell_0}{\partial \delta_0}$$

Ogni gradiente è un solo prodotto matrice-vettore, perché $\boldsymbol{\lambda}_{j+1}$ ha già accumulato tutta la sensibilità futura.

Il nome "adjoint" viene dall'analisi funzionale: $(\cdot)^\top$ è l'operatore aggiunto (adjoint) della Jacobiana. In italiano si dice anche "co-stato" perché in meccanica classica questo vettore gioca il ruolo delle variabili coniugate (i momenti) nello spazio delle fasi — evolve all'indietro nel tempo con la dinamica trasposta, proprio come il momento coniugato evolve con l'Hamiltoniana trasposta.
%%%%%%%%%%%%%%%%%%%%%%%%%%
\paragraph{Propagation through RK4.}
A subtlety shared with the linearisation in Section~\ref{sec:jacobian} is that
each step $\Psi_j$ is itself a four-stage composition (the RK4 sub-steps
$k_1, \ldots, k_4$), each of which calls LDNet.
Reverse-mode AD treats each sub-step as an additional block in the primitive
chain, so \eqref{eq:adjoint_rec} propagates \emph{exactly} through all four
stages and all neural-network activations, accumulating no truncation error.
The gradient returned by \texttt{GradientTape} therefore reflects the true
sensitivity of the discrete-time cost to the control sequence, not a
finite-difference approximation.

\paragraph{Why Adam instead of a QP solver.}
In the LQR case the cost is quadratic in the state and the dynamics are linear,
so the optimisation has a unique closed-form solution (the DARE
\eqref{eq:dare}).
For MPC with a nonlinear plant the cost \eqref{eq:mpc_cost} is generally
\emph{nonconvex} in $\mathbf{u}$: the aerodynamic output $C_L$ is a nonlinear
function of the state trajectory, which itself depends nonlinearly on
$\mathbf{u}$ through $\mathcal{R}$.
A QP solver would require a convex quadratic approximation (e.g.\ sequential
quadratic programming), introducing an inner loop and potentially large
linearisation errors when $z$ deviates significantly from trim.

Adam \citep{kingma2015adam} avoids this approximation by operating directly on
the nonconvex landscape via gradient descent with adaptive per-element learning
rates.
The first and second moment estimates $(m_t, v_t)$ implicitly rescale the
gradient by an approximation of the inverse diagonal Hessian, providing
faster convergence than vanilla gradient descent without requiring an explicit
Hessian computation:
\begin{align}
  m_{t+1} &= \beta_1 m_t + (1-\beta_1)\,\nabla_{\mathbf{u}} J_N, \label{eq:adam_m}\\
  v_{t+1} &= \beta_2 v_t + (1-\beta_2)\,(\nabla_{\mathbf{u}} J_N)^2, \label{eq:adam_v}\\
  \mathbf{u}_{t+1} &= \mathbf{u}_t - \eta\,\frac{\hat{m}_{t+1}}{\sqrt{\hat{v}_{t+1}} + \epsilon},
  \label{eq:adam_update}
\end{align}
where $\hat{m}$, $\hat{v}$ are bias-corrected estimates, $\eta = 0.5$ is the
learning rate, $\beta_1 = 0.9$, $\beta_2 = 0.999$, and $\epsilon = 10^{-8}$.
The box constraint $|\delta_j| \le \delta_{\max}$ is enforced by clipping
$\mathbf{u}$ before each forward pass:
$\mathbf{u}_{\rm cl} = \operatorname{clip}(\mathbf{u}, -\delta_{\max}, +\delta_{\max})$.
Because clipping is a piecewise-linear operation, TensorFlow propagates
gradients through it correctly (zero gradient at saturated elements), so no
constraint-handling logic is required beyond the clip.

\paragraph{Warm-starting.}
\label{sec:warm_start}
At timestep $k+1$ the optimal sequence $\mathbf{u}^*_k = [\delta_k^*, \ldots, \delta_{k+N-1}^*]^\top$
from the previous solve provides a good initial guess for the new problem.
The sequence is shifted by one:
\begin{equation}
  \mathbf{u}_0^{(k+1)} = [\delta_{k+1}^*, \delta_{k+2}^*, \ldots, \delta_{k+N-1}^*, 0]^\top,
  \label{eq:warm_start}
\end{equation}
appending a zero for the new final step.
This warm-start exploits the temporal continuity of the optimal policy: the
shifted sequence is already near-feasible for the new initial condition, so
Adam converges in fewer iterations than a cold start.
In practice $n_{\rm iter} = 15$ Adam steps from the warm-start reduce the cost
sufficiently for the solve to qualify as real-time.

\paragraph{Convergence and real-time feasibility.}
The fixed iteration budget $n_{\rm iter} = 15$ does \emph{not} guarantee
convergence to the true minimum of \eqref{eq:mpc_cost}; it trades solution
quality for computational speed.
Three factors mitigate the quality loss: (i) warm-starting means the initial
cost is already low, so few iterations are needed; (ii) the receding-horizon
re-planning at the next timestep corrects any residual suboptimality; and
(iii) the cost landscape is mildly nonconvex for small gust amplitudes, so the
warm-started Adam iterate is unlikely to be trapped in a poor local minimum.
The sensitivity of GLA performance to $n_{\rm iter}$ is explored in
Section~\ref{sec:mpc_sensitivity}.

The entire solve (rollout + backward pass + $n_{\rm iter}$ Adam steps) is
compiled as a single \texttt{@tf.function} to eliminate Python overhead.
Wall-clock times on a modern CPU are reported alongside the simulation results
in Section~\ref{sec:mpc_results}.

% ================================================================
\subsubsection{Simulation Results: MPC vs.\ LQR}
\label{sec:mpc_results}
% ================================================================

\paragraph{Setup.}
Two controllers are compared under the same 1-cosine gust scenario used for
the LQR evaluation:
\begin{equation}
  W(t) = \frac{W_0}{2}\Bigl(1 - \cos\tfrac{2\pi t}{T_g}\Bigr),
  \quad 0 \le t \le T_g,
\end{equation}
with $W_0 = 60\,\text{m/s}$, $T_g = 1\,\text{s}$, $U_\infty = 75\,\text{m/s}$.
\begin{enumerate}
  \item \textbf{LQR + EKF} --- best-performing configuration from
        Section~\ref{sec:results}: fixed offline gain with EKF observer.
        Taken directly from the previous chapter as the baseline.
  \item \textbf{MPC + EKF} --- MPC solver (Section~\ref{sec:mpc_solver})
        with the same EKF observer, horizon $N=10$, weights
        \eqref{eq:mpc_weights_state}--\eqref{eq:mpc_weights_ctrl},
        $n_{\rm iter}=15$ Adam steps, warm-started.
\end{enumerate}
Both controllers use the same observer (EKF, Section~\ref{sec:ekf}) to isolate
the effect of the control law on closed-loop performance.

\paragraph{Closed-loop performance.}
Figure~\ref{fig:mpc_performance} shows heave acceleration $\ddot{h}$, pitch
angle $\alpha$, flap deflection $\delta$, and flap rate $\dot{\delta}$ for all
four configurations.

\begin{figure}[htbp]
  \centering
  \includegraphics[width=\textwidth]{Chapters/Chapter3/Images/Section3/mpc_fig1_performance.png}
  \caption{Closed-loop performance comparison.
           Heave acceleration $\ddot{h}$, pitch angle $\alpha$, flap deflection
           $\delta$ ($+$\,=\,flap down), and flap rate $\dot{\delta}$.
           Blue: open-loop baseline. Steelblue: MPC oracle (true state + gust
           preview). Cyan: LQR+EKF. Red: MPC+EKF.}
  \label{fig:mpc_performance}
\end{figure}

\paragraph{Aerodynamic coefficients.}


\begin{figure}[htbp]
  \centering
  \includegraphics[width=\textwidth]{Chapters/Chapter3/Images/Section3/mpc_fig2_aero.png}
  \caption{Aerodynamic coefficient time histories.
           Blue: open-loop baseline. Steelblue: MPC oracle. Cyan: LQR+EKF.
           Red: MPC+EKF.}
  \label{fig:mpc_aero}
\end{figure}

\paragraph{Computational cost.}

\begin{figure}[htbp]
  \centering
  \includegraphics[width=\textwidth]{Chapters/Chapter3/Images/Section3/mpc_fig3_solve_time.png}
  \caption{Per-timestep MPC solve time (wall-clock, CPU).
           The first call includes \texttt{@tf.function} JIT compilation overhead.}
  \label{fig:mpc_solve_time}
\end{figure}

% ================================================================
\subsubsection{Sensitivity Analysis}
\label{sec:mpc_sensitivity}
% ================================================================

The MPC design involves three parameters with no counterpart in the LQR:
horizon length $N$, Adam learning rate $\eta$, and iteration count
$n_{\rm iter}$.
Their effects on GLA performance and computational cost are summarised below.

\paragraph{Horizon $N$.}


\begin{figure}[htbp]
  \centering
  \includegraphics[width=\textwidth]{Chapters/Chapter3/Images/Section3/mpc_fig4_horizon.png}
  \caption{Peak heave acceleration and peak $C_L$ reduction vs.\ horizon $N$.
           Solid lines: MPC+EKF. Dashed: MPC oracle. Dotted: LQR+EKF (horizon independent).}
  \label{fig:mpc_horizon}
\end{figure}

\paragraph{Learning rate $\eta$ and iteration count $n_{\rm iter}$.}


\begin{figure}[htbp]
  \centering
  \includegraphics[width=\textwidth]{Chapters/Chapter3/Images/Section3/mpc_fig5_adam.png}
  \caption{Cost $J_N$ vs.\ Adam iteration for several learning-rate / iteration-budget
           combinations. Warm-started from the previous step's solution.
           Baseline: $\eta = 0.5$, $n_{\rm iter} = 15$ (bold red).}
  \label{fig:mpc_adam}
\end{figure}

\paragraph{Cost weights.}


% ================================================================
\subsubsection{Conclusions}
\label{sec:mpc_conclusions}
% ================================================================

This chapter has presented a gradient-based MPC controller for gust load
alleviation on the LDNet aeroelastic model.
The key contributions relative to the LQR design are:

\begin{itemize}
  \item A formulation of finite-horizon nonlinear MPC as a sequence optimisation
        over the full nonlinear rollout of LDNet + structural RK4, with no
        linearisation required.
  \item An adjoint-based gradient derivation showing how reverse-mode AD
        through $N$ composed plant steps yields $\nabla_{\mathbf{u}} J_N$
        in $O(N)$ backward passes---the same computational primitive used
        for the LQR Jacobians in Section~\ref{sec:jacobian}, now applied online.
  \item An analysis of the Adam optimizer as the appropriate solver for this
        nonconvex, box-constrained problem: gradient clipping handles the
        actuator constraint without an inner loop, and the adaptive moment
        estimates provide implicit curvature rescaling without an explicit
        Hessian.
  \item A warm-starting strategy that exploits temporal continuity of the
        optimal sequence to reduce the required iteration budget to 15 Adam
        steps per timestep, achieving real-time feasibility on a standard CPU.
  \item Demonstration that MPC+EKF recovers a portion of the first $C_L$ peak
        that LQR+EKF cannot attenuate, at a computational overhead of
        approximately three orders of magnitude relative to the LQR law.
\end{itemize}

The primary limitation of the MPC implementation is its reliance on the linear
gust extrapolation \eqref{eq:W_extrap}: the remaining gap between MPC+EKF and
the oracle is dominated by gust estimation error, not by solver suboptimality.
Integrating a dedicated gust sensor (LiDAR upwash measurement or pitot array)
would close this gap significantly.
A secondary limitation is the fixed iteration budget: the Adam solver is not
guaranteed to converge to the global minimum of \eqref{eq:mpc_cost}, and for
very aggressive gust scenarios the warm-started iterate may enter a region
where the cost landscape is strongly nonconvex.
Increasing $n_{\rm iter}$ or adding a second-order correction step would
improve robustness at the cost of higher latency.

Both limitations point to the same practical conclusion: for a real-time
implementation with limited sensor suite and modest computational budget, the
LQR provides a compelling baseline that is within a few percent of MPC
performance during the post-gust oscillation phase.
MPC justifies its additional complexity primarily through first-peak attenuation,
which becomes material only when a sufficiently accurate forward gust estimate
is available.